% ============================================================================
\section{Controlling MountainCar}
\label{sec:mountaincar}
% Source: assignments/04_mountaincar/index.html
% ============================================================================

% ----------------------------------------------------------------------------
\subsection{Value iteration}

$\text{States}=\{-2,-1,0,+1,+2\}$, start state $0$, end states $\pm2$, $\text{Actions}=\{a_1,a_2\}$, $\gamma=1$.
For $i\in\{-1,0,1\}$: $T(i-1\mid i,a_1)=0.8$, $T(i+1\mid i,a_1)=0.2$; $T(i-1\mid i,a_2)=0.7$, $T(i+1\mid i,a_2)=0.3$.
$\text{Reward}(s,a,s')=10$ if $s'=-2$; $50$ if $s'=+2$; $-5$ otherwise.

\problempart{1a}{Two iterations}
\question{Give $V^\star_i(s)$ for every $s$ after iterations $i=1,2$ (with $V^\star_0\equiv0$, and $V^\star_i=0$ at end states).}
\expect{All 10 values.}
\begin{solution}
Each iteration applies $V^\star_i(s)=\max_a\sum_{s'}T(s'\mid s,a)\bigl[\text{Reward}(s,a,s')+V^\star_{i-1}(s')\bigr]$
($\gamma=1$) to the non-end states.

\textbf{Iteration 1} ($V^\star_0\equiv0$):
\begin{itemize}[nosep]
  \item $s=-1$: $Q(a_1)=0.8(10)+0.2(-5)=7$, \ $Q(a_2)=0.7(10)+0.3(-5)=5.5$ \ $\Rightarrow V^\star_1(-1)=7$
  \item $s=0$: $Q(a_1)=0.8(-5)+0.2(-5)=-5$, \ $Q(a_2)=0.7(-5)+0.3(-5)=-5$ \ $\Rightarrow V^\star_1(0)=-5$
  \item $s=+1$: $Q(a_1)=0.8(-5)+0.2(50)=6$, \ $Q(a_2)=0.7(-5)+0.3(50)=11.5$ \ $\Rightarrow V^\star_1(1)=11.5$
\end{itemize}
\textbf{Iteration 2} (using $V^\star_1$):
\begin{itemize}[nosep]
  \item $s=-1$: $Q(a_1)=0.8(10+0)+0.2(-5-5)=6$, \ $Q(a_2)=0.7(10)+0.3(-10)=4$ \ $\Rightarrow V^\star_2(-1)=6$
  \item $s=0$: $Q(a_1)=0.8(-5+7)+0.2(-5+11.5)=2.9$, \ $Q(a_2)=0.7(2)+0.3(6.5)=3.35$ \ $\Rightarrow V^\star_2(0)=3.35$
  \item $s=+1$: $Q(a_1)=0.8(-5-5)+0.2(50+0)=2$, \ $Q(a_2)=0.7(-10)+0.3(50)=8$ \ $\Rightarrow V^\star_2(1)=8$
\end{itemize}
\begin{center}
\begin{tabular}{lccccc}
\toprule
 & $s=-2$ & $s=-1$ & $s=0$ & $s=+1$ & $s=+2$\\
\midrule
$V^\star_1$ & 0 & 7 & $-5$ & 11.5 & 0\\
$V^\star_2$ & 0 & 6 & 3.35 & 8 & 0\\
\bottomrule
\end{tabular}
\end{center}
\end{solution}

\problempart{1b}{Optimal policy}
\question{Using $V^\star_2$, what is $\pi^\star(s)$ for each non-end state?}
\expect{$\pi^\star(s)$ for $s\in\{-1,0,+1\}$.}
\begin{solution}
Taking $\arg\max_a Q(s,a)$ with $Q$ computed from $V^\star_2$:
\begin{itemize}[nosep]
  \item $s=-1$: $Q(a_1)=0.8(10)+0.2(-5+3.35)=7.67$, \ $Q(a_2)=0.7(10)+0.3(-1.65)=6.505$ \ $\Rightarrow\pi^\star(-1)=a_1$
  \item $s=0$: $Q(a_1)=0.8(-5+6)+0.2(-5+8)=1.4$, \ $Q(a_2)=0.7(1)+0.3(3)=1.6$ \ $\Rightarrow\pi^\star(0)=a_2$
  \item $s=+1$: $Q(a_1)=0.8(-1.65)+0.2(50)=8.68$, \ $Q(a_2)=0.7(-1.65)+0.3(50)=13.845$ \ $\Rightarrow\pi^\star(1)=a_2$
\end{itemize}
So $\pi^\star(-1)=a_1$, $\pi^\star(0)=a_2$, $\pi^\star(+1)=a_2$. Near the left end the sure reward of $10$ is worth
taking, while from $0$ and $+1$ it pays to gamble on reaching $+50$. This matches the policy our
\texttt{value\_iteration} code returns for this MDP (grader test \texttt{3a-0-basic}).
\end{solution}

% ----------------------------------------------------------------------------
\subsection{Transforming MDPs}

\problempart{2a}{Removing the discount}
\question{Given an MDP with discount $\gamma<1$ and a solver that only handles $\gamma'=1$, define a new MDP with
$\text{States}'=\text{States}\cup\{o\}$, the same actions, and $\gamma'=1$. Specify $T'(s'\mid s,a)$ and
$\text{Reward}'(s,a,s')$ so that $V_{\text{opt}}(s)$ is unchanged for every $s\in\text{States}$.}
\expect{$T'$ and $\text{Reward}'$ as formulas, followed by a short verification that the optimal-value
recurrences coincide. Use the notation of the question.}
\begin{solution}
Interpret the discount as a probability $1-\gamma$ of the episode ending at every step: make $o$ an absorbing end
state ($\text{IsEnd}(o)$, so $V'_{\text{opt}}(o)=0$) and, for $s\in\text{States}$ (assuming $0<\gamma<1$),
\[
T'(s'\mid s,a)=
\begin{cases}\gamma\,T(s'\mid s,a) & s'\in\text{States}\\ 1-\gamma & s'=o\end{cases}
\]
\[
\text{Reward}'(s,a,s')=
\begin{cases}\dfrac{\text{Reward}(s,a,s')}{\gamma} & s'\in\text{States}\\[4pt] 0 & s'=o.\end{cases}
\]
$T'$ is a valid distribution because $\gamma\sum_{s'}T(s'\mid s,a)+(1-\gamma)=1$.
\emph{Verification.} For any non-end $s\in\text{States}$, the optimal values of the new MDP satisfy
\begin{align*}
V'_{\text{opt}}(s)
&=\max_{a}\Bigl[\sum_{s'\in\text{States}}\gamma\,T(s'\mid s,a)\Bigl(\tfrac{\text{Reward}(s,a,s')}{\gamma}+V'_{\text{opt}}(s')\Bigr)
  +(1-\gamma)\bigl(0+V'_{\text{opt}}(o)\bigr)\Bigr]\\
&=\max_{a}\sum_{s'\in\text{States}}T(s'\mid s,a)\bigl[\text{Reward}(s,a,s')+\gamma\,V'_{\text{opt}}(s')\bigr],
\end{align*}
using $V'_{\text{opt}}(o)=0$. This is exactly the Bellman optimality equation of the original MDP, and end states
of the original MDP keep value $0$ in both. For $\gamma<1$ that equation has a unique solution (the Bellman
operator is a contraction), so $V'_{\text{opt}}(s)=V_{\text{opt}}(s)$ for all $s\in\text{States}$, and the optimal
policies coincide. (Equivalently, one may keep $\text{Reward}'=\text{Reward}$ on $\text{States}$ and set
$\text{Reward}'(s,a,o)=\sum_{s'}T(s'\mid s,a)\,\text{Reward}(s,a,s')$; this avoids dividing by $\gamma$.)
\end{solution}

% ----------------------------------------------------------------------------
\subsection{Value iteration on Mountain Car}
\codepart{4.3a--4.3b}
\texttt{value\_iteration} and \texttt{ModelBasedMonteCarlo}. All 11 points of grader tests \texttt{3a}/\texttt{3b} pass,
and the hidden speed test ($n=500$) runs in 2.6\,s (limit 14\,s).

\paragraph{How the experiments were run.} \texttt{scripts/mountaincar\_experiments.py} uses exactly the
configurations of \texttt{train.py}: a 20$\times$20 discretisation for value iteration and tabular Q-learning, and
Fourier features ($m=5$, scale $[1,15]$) for function approximation, with discount $0.999$ and episodes of at most 1000
steps. Each configuration is trained for 1000 episodes and then evaluated for 500 greedy episodes. We seed each of
the three trials so that the figures and numbers below are reproducible, and we count an evaluation episode as
successful when it ends in the goal state rather than at the time limit.

\problempart{3c}{Model-based value iteration: results}
\question{Run \texttt{python train.py -{}-agent value-iteration} (three trials). Comment on the reward plots, and
discuss when model-based value iteration can perform poorly in general.}
\expect{The plots, plus 2--3 sentences describing them and when the method may fail.}
\begin{solution}
\begin{figure}[H]\centering
  \includegraphics[width=0.32\linewidth]{04_mountaincar/mcvi_1.png}\hfill
  \includegraphics[width=0.32\linewidth]{04_mountaincar/mcvi_2.png}\hfill
  \includegraphics[width=0.32\linewidth]{04_mountaincar/mcvi_3.png}
  \caption{Model-based value iteration, trials 1--3 (blue: training episodes, green: greedy evaluation episodes;
  lines are 30-episode moving averages).}
\end{figure}
\begin{center}\small
\begin{tabular}{cccc}
\toprule
trial & train mean (last 100 ep.) & eval mean $\pm$ std & eval episodes reaching the goal\\
\midrule
1 & $-238.5$ & $-128.8\pm76.1$ & 495/500\\
2 & $-240.3$ & $-147.8\pm71.5$ & 495/500\\
3 & $-232.1$ & $-130.1\pm76.9$ & 495/500\\
\bottomrule
\end{tabular}
\end{center}
All three trials show the same staircase. For the first $\approx100$ episodes the agent acts randomly and never reaches
the flag (reward $\approx-750$, the 1000-step time limit). The reward then jumps up twice, at episodes $\approx100$
and $\approx200$: these are exactly the moments value iteration runs (every $10^5$ steps, and each failed episode
lasts 1000 steps) and installs a new policy built from the estimated model. After the second solve, training rewards
plateau around $-235$. They stay noisy because the agent still explores with probability $0.5$, and the policy is
fixed between solves. The greedy evaluation is much better ($\approx-130$, goal reached in 495/500 episodes), with
rewards falling into a few bands that correspond to the different discretised start positions.

Model-based value iteration can perform poorly when the model it estimates is wrong or incomplete. That happens when
the state space is large or continuous (many state--action pairs are visited rarely or never, so $\hat T,\hat R$ are
estimated from very few samples), when the discretisation hides relevant information so that the ``discrete MDP'' is
not really Markov, when exploration does not reach important regions (states never visited get no policy), and when
the environment is non-stationary so that old counts mislead the model. It is also expensive: every re-solve costs a
full value iteration over all states.
\end{solution}

% ----------------------------------------------------------------------------
\subsection{Q-learning Mountain Car}
\codepart{4.4a--4.4c}
\texttt{TabularQLearning}, \texttt{fourier\_feature\_extractor}, \texttt{FunctionApproxQLearning}. All 23 points of
grader tests \texttt{4a}--\texttt{4c} pass.

\problempart{4d}{Tabular vs.\ function approximation}
\question{Run \texttt{train.py -{}-agent tabular} and \texttt{-{}-agent function-approximation}. Compare the plots.
Tabular Q-learning is expected to do better on this task. Why might that be?}
\expect{The plots, plus 2--3 sentences comparing the two and giving possible reasons.}
\begin{solution}
\begin{figure}[H]\centering
  \includegraphics[width=0.32\linewidth]{04_mountaincar/tabular_1.png}\hfill
  \includegraphics[width=0.32\linewidth]{04_mountaincar/tabular_2.png}\hfill
  \includegraphics[width=0.32\linewidth]{04_mountaincar/tabular_3.png}\\[4pt]
  \includegraphics[width=0.32\linewidth]{04_mountaincar/fapprox_1.png}\hfill
  \includegraphics[width=0.32\linewidth]{04_mountaincar/fapprox_2.png}\hfill
  \includegraphics[width=0.32\linewidth]{04_mountaincar/fapprox_3.png}
  \caption{Top: tabular Q-learning; bottom: function-approximation Q-learning (Fourier features); trials 1--3.}
\end{figure}
\begin{center}\small
\begin{tabular}{lcccc}
\toprule
agent & trial & train mean (last 100) & eval mean $\pm$ std & eval goal reached\\
\midrule
tabular & 1 & $-228.4$ & $-193.8\pm37.7$ & 500/500\\
tabular & 2 & $-240.7$ & $-268.8\pm161.0$ & 446/500\\
tabular & 3 & $-198.8$ & $-159.2\pm3.2$ & 500/500\\
\midrule
function approx. & 1 & $-375.5$ & $-367.3\pm50.3$ & 500/500\\
function approx. & 2 & $-694.4$ & $-697.6\pm18.9$ & 0/500\\
function approx. & 3 & $-321.8$ & $-360.3\pm132.4$ & 484/500\\
\bottomrule
\end{tabular}
\end{center}
Returns are discounted ($\gamma=0.999$) and all runs are seeded with the trial number
(\texttt{scripts/mountaincar\_experiments.py}). Tabular Q-learning is clearly better here. Its evaluation returns
lie between $-159$ and $-269$, and two trials reach the goal in every evaluation episode. Function approximation
reaches $-360$ at best and is far less reliable: trial~2 never learns to escape the valley at all. Its curve stays
flat near $-700$, which is the discounted return of a 1000-step timeout. Two reasons fit what we see:
\begin{itemize}
  \item \textbf{Approximation error.} The 36 Fourier basis functions (\texttt{max\_coeff}$=5$) can only represent
  a smooth Q-function. The true one has a sharp edge between states that have enough energy to crest the hill and
  states that do not. The $20\times20$ table has 400 independent cells and can represent that edge, and with only a
  2-D state space, $\epsilon$-greedy exploration visits nearly all of them in 1000 episodes.
  \item \textbf{Interference.} Every Fourier feature is global, so each semi-gradient update changes Q-values
  everywhere, and learning in one region can undo learning in another. Q-learning with function approximation has
  no convergence guarantee, which explains the large differences between trials.
\end{itemize}
\end{solution}

\problempart{4e}{When function approximation wins}
\question{Why could function-approximation Q-learning beat tabular Q-learning in other environments, e.g.\ with
limited exploration, very high-dimensional state spaces, or tight memory?}
\expect{2--3 sentences.}
\begin{solution}
A tabular agent learns nothing about a state--action pair until it has visited it, and it needs one entry per
discretised state. Both grow exponentially with the state dimension ($20^d$ cells for $d$ variables), so in
high-dimensional spaces or with a limited exploration budget most cells are never visited and the table
cannot even be stored. Function approximation \emph{generalises}: an update from one state changes the Q-values of
all similar states through the shared weights, so a good policy can be learned from far fewer samples. Its memory
is fixed by the number of features (here $36\times3$ weights instead of $20^2\times3$ table entries), not by the
size of the state space, and it works directly on continuous states without discretisation error.
\end{solution}

% ----------------------------------------------------------------------------
\subsection{Safe exploration}

\problempart{5a}{\texttt{max\_speed} in Gym}
\question{Explain how \texttt{self.max\_speed} is used in Gym's \texttt{mountain\_car.py}.}
\expect{One sentence.}
\begin{solution}
After every velocity update the environment clips the velocity to $[-\texttt{max\_speed},\,\texttt{max\_speed}]$
(\texttt{np.clip(velocity, -self.max\_speed, self.max\_speed)}), so the car can never move faster than
\texttt{max\_speed} in either direction. The value also sets the velocity bounds of the observation space, which our
reward function uses (the velocity term is $-(\texttt{max\_speed}-|v|)$).
\end{solution}

\problempart{5b}{Removing the constraint}
\question{Run \texttt{train.py -{}-agent function-approximation -{}-max\_speed=100000}. Why does the behaviour of the
learned policy barely change compared with $\texttt{max\_speed}=0.07$?}
\expect{1--2 sentences.}
\begin{solution}
Removing the limit does not change the physics. With \texttt{force}$=0.001$ and \texttt{gravity}$=0.0025$, the car
never comes close to $0.07$. We ran the environment with \texttt{max\_speed}$=10^5$ (no clipping) and recorded the
largest speed reached:
\begin{itemize}
  \item an energy-pumping policy (push in the direction of motion), 200 episodes: $|v|_{\max}=0.062$;
  \item a random policy, 50 episodes: $|v|_{\max}=0.057$.
\end{itemize}
So the clip never binds, and the transition dynamics, and therefore the optimal policy, stay the same.

What does change is the reward. The shaping term $-(\texttt{max\_speed}-|v|)$ becomes an almost constant
$\approx-10^5$ per step. That rescales the returns to the order of $-10^7$ (trial~3's
$-63\,231\,126$ is exactly $10^5\sum_{t<1000}0.999^t$, a timeout in every episode), but it does not change which
action is best, because the only thing that matters is still reaching the goal in as few steps as possible. Two of
our three runs learned to reach the goal, in 417/500 and 499/500 evaluation episodes, and one failed. This is the
same pattern as the default-speed runs in 4d (trial~2 failed there too). The large reward scale makes
each TD update about $10^5$ times larger and swamps the small velocity bonus, so learning is somewhat noisier, but
the behaviour is qualitatively unchanged.
\begin{figure}[H]\centering
  \includegraphics[width=0.32\linewidth]{04_mountaincar/fapprox_maxspeed100000_1.png}\hfill
  \includegraphics[width=0.32\linewidth]{04_mountaincar/fapprox_maxspeed100000_2.png}\hfill
  \includegraphics[width=0.32\linewidth]{04_mountaincar/fapprox_maxspeed100000_3.png}
  \caption{Function-approximation Q-learning with \texttt{max\_speed}$=10^5$, trials 1--3.}
\end{figure}
\end{solution}

\problempart{5c}{Constrained Q-learning}
\question{After implementing \texttt{ConstrainedQLearning}, so that only actions with
$\text{velocity}_{t+1}<\text{threshold}$ are allowed, with
$\text{velocity}_{t+1}=\text{velocity}_t+(\text{action}-1)\cdot\text{force}-\cos(3\,\text{position}_t)\cdot\text{gravity}$,
run \texttt{grader.py 5c-helper}. Why do the policies for $\texttt{max\_speed}=0.065$ and an unconstrained
\texttt{max\_speed} differ?}
\expect{1--2 sentences (only the written answer is graded). Code: notebook \S4.5c.}
\begin{solution}
In our run of \texttt{5c-helper}, the unconstrained agent ($\texttt{max\_speed}=10000$) used left/none/right
$136/152/142$ times in its final trajectory (430 steps), while the agent constrained to $0.065$ used them
$75/64/169$ times (308 steps). The two policies differ because the constraint changes the problem itself, not just
the learner. In every state where an action would push the predicted velocity to $0.065$ or above, that action is
removed, so near the top of a swing the constrained agent cannot pick the same greedy action as the unconstrained
one. It is also only ever trained on, and evaluates Q-values over, the restricted action set, so it learns different
Q-values and settles on a different, more right-heavy way of building momentum that stays within the speed limit.
The counts change from run to run because the agents are trained online for only 100 episodes, but the
constrained action sets, and hence the policies, are genuinely different.
\end{solution}

\problempart{5d}{Designing safe MDPs for autonomous driving}
\question{An autonomous car is trained with RL on public residential roads (pedestrians, vehicles, traffic signals;
camera sensing). Design \textbf{MDP A} with an exploration policy that could cause harm, and \textbf{MDP B}, a
modification that reduces or avoids that harm.}
\expect{For each MDP: $\text{States}$, $\text{Actions}(s)$, $\text{Reward}(s,a,s')$, an exploration policy, and
a 1--2 sentence ethical explanation.}
\begin{solution}
\textbf{MDP A (harmful).}
\begin{itemize}[nosep]
  \item \emph{States:} $s=(x, v, \theta, \text{lane}, \text{light}, P, O)$: the car's position, speed and heading, current
        lane, traffic-light state ahead, and the detected positions and velocities of pedestrians $P$ and other
        vehicles $O$ from the cameras.
  \item \emph{Actions:} $\text{Actions}(s)=\{\text{accelerate}, \text{brake}, \text{steer left}, \text{steer right},
        \text{keep}\}$ in every state, with no restrictions.
  \item \emph{Reward:} $\text{Reward}(s,a,s')=-1$ per time step, $+100$ on reaching the destination, and $-10$ for a
        collision; nothing for speeding, red lights or pedestrians.
  \item \emph{Exploration:} $\epsilon$-greedy with $\epsilon=0.3$: in 30\,\% of time steps a uniformly random action
        is executed on the real road.
\end{itemize}
\emph{Why harmful:} random actions on a public road mean random braking, swerving into other lanes or accelerating
through crossings, and the reward makes a red light or a pedestrian in the way merely a small cost compared with
arriving faster. The people who bear the risk of this exploration (pedestrians and other drivers) never consented to
it.

\textbf{MDP B (safer).}
\begin{itemize}[nosep]
  \item \emph{States:} the same as in MDP A, plus the local speed limit $v_{\max}$ and the minimum distance $d_{\min}$
        to any pedestrian or vehicle.
  \item \emph{Actions:} $\text{Actions}(s)$ contains only actions that a rule-based safety filter certifies as safe
        for the next step. The predicted speed stays below $v_{\max}$, the predicted distance to every pedestrian
        and vehicle stays above a safety margin, and the car may not enter an intersection on red. If no action is
        certified, only ``brake'' is allowed (constrained RL, as in 5c).
  \item \emph{Reward:} $-1$ per time step, $+100$ at the destination, $-1000$ for any collision or traffic violation, and
        a smooth penalty $-c\cdot\max(0, d_{\text{safe}}-d_{\min})$ for getting close to a pedestrian.
  \item \emph{Exploration:} explore only in simulation and on closed test tracks. On public roads, run greedily with
        a human safety driver able to take over at any time; any remaining exploration ($\epsilon\le0.05$) samples
        only among the certified actions.
\end{itemize}
\emph{Why safer:} harm is prevented by construction rather than merely discouraged, because unsafe actions are never
available, and the risky trial-and-error happens where nobody can be hurt. The reward now encodes the values society
expects (laws, distance to vulnerable road users), so the agent does not learn that violations are cheap.
\end{solution}
